% Generated draft: requires mathematical and source review.
\documentclass[11pt]{article}
\usepackage[margin=25mm]{geometry}
\usepackage{fontspec}
\setmainfont{Latin Modern Roman}
\usepackage{amsmath,amssymb,mathtools,unicode-math}
\setmathfont{Latin Modern Math}
\usepackage{enumitem,needspace,xurl,hyperref}
\setlength{\parindent}{0pt}
\setlength{\parskip}{5pt}
\setlength{\emergencystretch}{3em}
\newcommand{\N}{\mathbb{N}}
\newcommand{\Z}{\mathbb{Z}}
\newcommand{\Q}{\mathbb{Q}}
\newcommand{\R}{\mathbb{R}}
\newcommand{\C}{\mathbb{C}}
\newcommand{\F}{\mathbb{F}}
\newcommand{\A}{\mathbb{A}}
\newcommand{\PP}{\mathbb{P}}
\newcommand{\E}{\mathbb{E}}
\newcommand{\eps}{\varepsilon}
\newcommand{\dd}{\,\mathrm{d}}
\begin{document}
\section*{UnsolvedMath ID: 701678}
\textbf{Source problem code:} catalogue entry
\subsection{Definitions and mathematical statement}
For positive integers \(l,m,n\), define
\[
E(l,m,n)=\langle a,b\mid a^l=1,\ (ab)^m=b^n\rangle.
\]
The supplied record does not explicitly state the domain of the parameters. The formulation below uses the usual intended interpretation \(l,m,n\in\Z_{>0}\); this interpretation requires confirmation against the original source.

To specify the automaticity convention, let \(A\) be a finite, linearly ordered generating alphabet for a group \(G\), with evaluation map \(\pi:A^*\to G\), and assume that the image of \(A\) is closed under inverses. Order words first by length and then lexicographically. Let \(L\subseteq A^*\) contain the least word representing each element of \(G\) in this order. Write \(\operatorname{pad}(u,v)\) for the word of paired symbols obtained by padding the shorter of \(u,v\) on the right with a new symbol \(\diamond\). The group is \emph{shortlex automatic with respect to this ordered alphabet} if \(L\) is regular and, for every \(s\in A\cup\{\varepsilon\}\), the language
\[
\left\{\operatorname{pad}(u,v):u,v\in L,\ \pi(u)\pi(s)=\pi(v)\right\}
\]
is regular, where \(\varepsilon\) is the empty word and \(\pi(\varepsilon)=1\). Here regular means recognized by a finite automaton.

Under the convention that a group is shortlex automatic when such an ordered finite generating alphabet exists, the question is:
\[
\text{Does there exist }(l,m,n)\in\Z_{>0}^3
\text{ such that }\frac1l+\frac1m+\frac1n<1
\text{ and }E(l,m,n)\text{ is shortlex automatic?}
\]
The supplied record does not specify whether shortlex automaticity instead refers to the presentation alphabet \(\{a,a^{-1},b,b^{-1}\}\), or which ordering of that alphabet is intended. If that restricted meaning is intended, the generating alphabet and the quantifier over its orderings must be specified before the target is fully determined.

The record additionally reports that \(E(6,2,6)\) is not shortlex automatic. This is preserved as a supplied background assertion; its proof, reference, and generating-alphabet convention are not supplied. The registry labels the problem open; that label is supplied source metadata, not an independently verified mathematical assessment.
\subsection{Short English statement}
C. Chalk asks whether any parameters satisfying \(1/l+1/m+1/n<1\) yield a shortlex automatic group \(\langle a,b\mid a^l=1,\ (ab)^m=b^n\rangle\). Positive integer parameters are the provisional interpretation. The intended generating alphabet and ordering convention require confirmation. The source reports \(E(6,2,6)\) as a negative example.
\subsection{Sources}
\begin{itemize}
\item Local UnsolvedMath registry, permanent ID \(701678\), exact source problem code \texttt{catalogue entry}, title \emph{Kourovka Notebook Problem 20.20}. The supplied record labels its status \texttt{open} and gives upstream revision \texttt{d371ee917ac9647fbb06cd6fc529deec6f3dda82}.
\item Supplied source attribution: \emph{The Kourovka Notebook (Unsolved Problems in Group Theory)}, 20th issue (2022), problem 20.20; proposed by C. Chalk in 2022. The background also identifies the 21st issue edition (2026), \texttt{arXiv:1401.0300v48}. The collection description names E. I. Khukhro and V. D. Mazurov as editors.
\item The supplied statement reports that \(E(6,2,6)\) is not shortlex automatic, without a separate citation or proof. The named source and arXiv version have not been independently consulted for this draft.
\end{itemize}
\end{document}
